\chapter{Sökprotokoll för litteraturförankringen}
\label{app:sokprotokoll}

De pedagogiska designval som motiveras med litteratur i marginalnoterna
(\mintinline{latex}{\ltnote}) bygger på källor som sökts, lästs och
verifierats enligt principen \emph{finn $\to$ verifiera $\to$ registrera}:
varje referens har hittats via en dokumenterad sökning, dess fulltext har
lästs, och ett ordagrant citat som styrker påståendet har registrerats.
Per-referens\-provenienserna (\texttt{FOUND-VIA}, \texttt{QUOTE},
\texttt{VERIFIED}) står som kommentarer ovanför respektive post i
\texttt{ltnotes.bib}.  Detta appendix sammanfattar sökepisoderna.

Sökningarna gjordes med \texttt{scholar}-verktyget mot OpenAlex, Semantic
Scholar och DBLP; fulltexter hämtades öppet och citerades med
\texttt{scholar pdf quote} eller lästes direkt.  Datum: 2026-07-20.

\begin{description}
  \item[Strukturerat styrflöde (sekvens, val, upprepning)]
    Fråga: \enquote{Flow diagrams Turing machines languages two formation
    rules}.  Behållen: \textcite{bohm1966} (det strukturerade
    programmeringsteoremet), verifierad mot den öppna avskriften
    \url{https://www.cs.unibo.it/~martini/PP/bohm-jac.pdf}; citatet om
    \enquote{composition and iteration} och de tre bas\-diagram\-typerna
    lästes i fulltext.

  \item[Sekvens/exekveringsordning som kritisk aspekt]
    Fråga: \enquote{Pea language-independent conceptual bugs novice
    programming}.  Behållen: \textcite{pea1986} (parallellism-buggen och
    \enquote{superbuggen}), verifierad i öppen avskrift; citatet om att
    rader antas vara aktiva samtidigt lästes i fulltext.  Avskriften har
    OCR-artefakter, som rensats till läsbar ordalydelse i proveniensen.

  \item[CS1:s grundkonstruktioner och sekventiell exekvering]
    Fråga: \enquote{Qian Lehman students misconceptions introductory
    programming literature review}.  Behållen: \textcite{qian2017}
    (översikt som räknar upp variabler, villkor, slingor, sekventiell
    exekvering och funktioner), verifierad genom \texttt{pdftotext} av
    ACM-pdf:en; meningen \enquote{program instructions are executed
    sequentially} lästes direkt.

  \item[Variationsteorins kontrastprincip]
    \textcite{Marton2015} används kursövergripande i marginalnoterna och
    verifierades i variationsteori-underlaget (se \texttt{ltnotes.bib}).
\end{description}

\paragraph{Sökt men ej använt.}
En bredare sökning (för att pröva om listan är fullständig) tog upp
du~Boulays begrepp \enquote{notional machine} (1986), Sorvas arbete om
skiftet från statisk till dynamisk exekvering (2012, 2013) och Eckerdals
variationsteoretiska analyser av att lära sig programmera (2009).  Dessa
stärker slutsatsen att sekvens/exekveringsordning --- och mer allmänt
\enquote{hur datorn kör programmet} --- är kritiska aspekter.  De
\emph{citeras dock inte} här: fulltexterna kunde inte alla hämtas och
verifieras personligen vid detta tillfälle, och en overifierad referens
förs inte in.  Den notionella maskinen är den starkaste kandidaten till en
framtida utökning av lärandemålet, när källan väl lästs i original.
